
Para la implementación del gestor de monitoreo se seleccionó
\textbf{Ubuntu Server 26.04 LTS de 64 bits para arquitectura ARM64}.
Este sistema operativo será instalado en una Raspberry Pi 5 con 8~GB de
memoria RAM y una unidad SSD externa de 120~GB.

Ubuntu Server es una distribución de GNU/Linux desarrollada por Canonical y
orientada al funcionamiento de servidores. A diferencia de Ubuntu Desktop,
su instalación predeterminada no incluye un entorno gráfico completo, por
lo que consume menos memoria RAM, capacidad de almacenamiento y recursos del
procesador. Esta característica permite que una mayor cantidad de recursos
del equipo permanezca disponible para Netdata y para los demás servicios de
monitoreo.

La versión seleccionada pertenece a la categoría LTS, correspondiente a
\textit{Long Term Support} o soporte de largo plazo. Las versiones LTS están
diseñadas para entornos que requieren estabilidad y mantenimiento durante
periodos prolongados, por lo que reciben actualizaciones de seguridad,
correcciones de errores y mantenimiento continuo. Esto reduce la necesidad
de realizar cambios frecuentes de versión y contribuye a mantener estable
la plataforma de monitoreo. \cite{ubuntu-server}

Ubuntu ofrece una imagen específica para Raspberry Pi en arquitectura ARM64,
la cual es compatible con el procesador de 64 bits de la Raspberry Pi 5
\cite{ubuntu-raspberry}. Esto permite aprovechar adecuadamente la capacidad
del procesador y los 8~GB de memoria RAM disponibles en el equipo.

Otro factor importante es la compatibilidad con Netdata. Ubuntu 26.04 y la
arquitectura AArch64 se encuentran entre las plataformas que cuentan con
soporte oficial por parte de Netdata \cite{netdata-plataformas}. Además,
Ubuntu utiliza paquetes de tipo \texttt{DEB}, el administrador de paquetes
\texttt{APT} y el administrador de servicios \texttt{systemd}, herramientas
compatibles con el proceso de instalación y ejecución del agente de Netdata.

Ubuntu Server también permite administrar el equipo de manera remota mediante
SSH. De esta forma, después de realizar la instalación inicial, no es
necesario mantener conectados permanentemente un monitor, teclado y ratón a
la Raspberry Pi. La configuración, actualización y supervisión del servidor
pueden realizarse desde otra computadora conectada a la misma red.

La utilización de una unidad SSD externa de 120~GB permitirá almacenar el
sistema operativo, los archivos de Netdata y el historial de métricas. En
comparación con una tarjeta microSD, el SSD proporciona una mayor velocidad
de lectura y escritura y resulta más apropiado para las operaciones frecuentes
de almacenamiento generadas por la base de datos interna de Netdata.

Por las razones anteriores, Ubuntu Server 26.04 LTS constituye una opción
adecuada para el proyecto, debido a que ofrece compatibilidad con la
Raspberry Pi 5, bajo consumo de recursos, estabilidad, administración remota,
soporte para ARM64 y compatibilidad oficial con Netdata.